\section{Introduction}\label{sec:intro}

\subsection{Background and Motivation}
In many Taiwanese cities, household waste is not left at the curb but handed directly to a collection truck that stops briefly at designated points according to a published schedule. Residents must therefore be physically present at a stop within a narrow time window. Taichung City publishes the positions of its garbage and recycling trucks as open data, with records containing route, vehicle, timestamp, location, coordinates, and speed fields refreshed approximately every ten minutes~\cite{taichung_fleet_api}. Such data make it possible to see \emph{where} a truck is, but they do not answer the question residents actually face: \emph{given where I am and how fast I walk, should I leave now, and for which stop?}

This question is becoming more consequential as the population ages. Taiwan crossed the threshold of a super-aged society, in which more than 20\% of the population is aged 65 or above, in 2025~\cite{ndc_population}, and walking speed varies substantially with age~\cite{knoblauch1996}. For an older resident, leaving too early means waiting at the curb, while leaving too late means missing the truck and storing waste for another day. A departure recommendation that is wrong in the optimistic direction is therefore more harmful than no recommendation at all.

\subsection{Problem Statement and Research Gap}
A departure decision combines three uncertain quantities: the truck's arrival time at a stop, the resident's walking time to that stop, and the freshness of the data on which both estimates depend. Existing tools address these separately. Municipal applications show static schedules or the latest vehicle icon but do not intersect vehicle arrival with pedestrian walking time. General navigation services estimate walking routes but have no model of collection-route logic or dwell behavior at stops. Large language models (LLMs) can converse naturally about the problem, but without grounded access to live data and spatial computation they may produce plausible yet fabricated times and distances~\cite{ji2023hallucination}.

Prior work on vehicle arrival prediction largely treats the task as improving point accuracy~\cite{chien2002,jeong2004}. With location updates arriving only every ten minutes, however, accuracy has a hard ceiling, and the more useful question becomes one of \emph{abstention}: under what data conditions is a system entitled to tell a user to leave, and when must it decline and fall back to less specific information? This framing follows the classical reject option in decision theory~\cite{chow1970} and selective prediction~\cite{geifman2017}, but to our knowledge it has not been applied to curbside collection or to tool-using LLM agents for municipal services.

\subsection{Research Contributions}
This paper presents \emph{GarbageGo}, a domain-specific, tool-grounded agent for departure decisions in curbside waste collection in Taichung City. Its contributions are as follows:
\begin{itemize}
\item \textbf{An advise-or-defer decision formulation.} We define departure feasibility through a time margin (slack) between predicted truck arrival, estimated walking time, and a safety buffer that widens with data latency, prediction error, and positioning uncertainty, together with explicit conditions under which the system must defer rather than advise.
\item \textbf{A tool-grounded agent architecture.} Domain knowledge, including schedules, pedestrian accessibility, and collection-operation characteristics, is encapsulated as whitelisted tools with JSON Schema-validated inputs and outputs. The LLM is restricted to intent recognition and explanation, and every number it reports must originate from a tool result.
\item \textbf{A safety-oriented evaluation protocol.} We specify a rolling-origin backtest that compares arrival-time predictors, treats the false-feasibility rate as the primary safety metric alongside recommendation coverage and fallback success, and includes at least thirty anomaly scenarios covering stale data, API failure, positioning error, and out-of-scope requests.
\end{itemize}
This paper reports the problem formulation, system design, and evaluation protocol. Empirical results will be reported after the backtest is completed; no experimental values are claimed here.

\subsection{Paper Organization}
Section~\ref{sec:related} reviews related work. Section~\ref{sec:method} presents the proposed scheme. Section~\ref{sec:experiments} describes the evaluation design, and Section~\ref{sec:conclusion} concludes the paper.
